\documentclass[11pt,a4paper]{article}
\usepackage[utf8]{inputenc}
\usepackage[T1]{fontenc}
\usepackage[french]{babel}
\usepackage[a4paper,margin=2.2cm]{geometry}
\usepackage{lmodern}
\usepackage{microtype}
\usepackage{xcolor}
\usepackage{tikz}
\usetikzlibrary{positioning,arrows.meta,calc,shapes.geometric}
\usepackage{listings}
\usepackage{enumitem}
\usepackage{hyperref}
\hypersetup{
  colorlinks=true, urlcolor=blue!60!black, linkcolor=blue!60!black,
  pdftitle={HearMyHands --- Flow des mod\`eles}
}

\definecolor{accent}{HTML}{4A90E2}
\definecolor{accent2}{HTML}{50E3C2}
\definecolor{warn}{HTML}{FF6B6B}
\definecolor{good}{HTML}{32B47D}
\definecolor{bg}{HTML}{F5F7FA}
\definecolor{code}{HTML}{2C3E50}

\lstset{
  basicstyle=\ttfamily\footnotesize\color{code},
  backgroundcolor=\color{bg},
  frame=single, framerule=0pt, framesep=6pt,
  breaklines=true, columns=fullflexible,
  keywordstyle=\color{accent}\bfseries,
  stringstyle=\color{good},
  commentstyle=\color{gray}\itshape,
}

\title{\textbf{HearMyHands} \\ \large Flow des donn\'ees entre les mod\`eles}
\author{\'Equipe HearMyHands --- Polytech Nantes}
\date{}

\begin{document}
\maketitle

\section{Architecture g\'en\'erale}

L'application repose sur \textbf{trois mod\`eles cha\^in\'es} c\^ot\'e backend
et un mod\`ele auxiliaire c\^ot\'e client en mode dynamique. Le navigateur
envoie ses frames webcam au serveur Flask via Socket.IO en binaire, et le
serveur r\'epond en JSON \`a chaque frame.

\begin{center}
\begin{tikzpicture}[
  font=\small,
  node distance=0.9cm and 1.2cm,
  every node/.style={align=center},
  box/.style={draw, rounded corners=3pt, minimum height=10mm, minimum width=22mm, fill=bg, drop shadow={opacity=0.15, shadow xshift=1pt, shadow yshift=-1pt}},
  modelbox/.style={box, fill=accent!15, draw=accent},
  clientbox/.style={box, fill=accent2!20, draw=accent2!80!black},
  arrow/.style={-{Stealth[length=2.2mm]}, line width=0.4mm},
]
  \node[clientbox] (browser) {Navigateur \\ \footnotesize webcam $\to$ JPEG 480px};
  \node[modelbox, right=2.6cm of browser] (vision) {\textbf{Vision} \\ \footnotesize CSPDarknet + heatmaps};
  \node[modelbox, below=of vision] (mp) {\textbf{MediaPipe Hands} \\ \footnotesize 21 landmarks $\times$ 2};
  \node[modelbox, right=1.6cm of vision] (mlp) {\textbf{MLP statique} \\ \footnotesize 42 $\to$ 30 $\to$ 23};
  \node[modelbox, right=1.6cm of mp] (gru) {\textbf{GRU Ocarina} \\ \footnotesize 45 frames $\times$ 42};
  \draw[arrow] (browser) -- node[above]{\footnotesize frame} (vision);
  \draw[arrow] (vision) -- node[right]{\footnotesize crop poignets} (mp);
  \draw[arrow] (mp) -| node[pos=0.25, above]{\footnotesize landmarks} (mlp);
  \draw[arrow] (mp) -- node[above]{\footnotesize buffer} (gru);
  \draw[arrow, dashed] (mlp.east) -- ++(0.6,0) |- node[pos=0.75, above right]{\footnotesize \texttt{\{letter, conf\}}} (browser.east);
  \draw[arrow, dashed] (gru.east) -- ++(0.4,0) |- node[pos=0.85, below right]{\footnotesize \texttt{\{sign, conf\}}} (browser.south east);
\end{tikzpicture}
\end{center}

\paragraph{Transport.} Frames JPEG ($\sim$480~px, qualit\'e 0{,}7) en binaire
sur Socket.IO. Pipeline avec backpressure : 4 frames en vol simultan\'ement
max, ack par s\'equence, application uniquement de la plus r\'ecente.

\section{Mode statique --- MLP par frame}

\textbf{Objectif} : reconna\^itre une lettre par frame, fluide en continu.
Utilis\'e pour les lettres statiques de l'alphabet LSF (toutes sauf J, P, X,
Z qui demandent un mouvement).

\subsection*{Flux serveur \`a chaque frame}

\begin{enumerate}[leftmargin=*, itemsep=2pt]
  \item D\'ecodage JPEG re\c{c}u.
  \item \textbf{Vision} : CSPDarknet53 + heatmaps + Soft-Argmax produisent
        9 keypoints (cou, \'epaules, coudes, poignets, hanches).
  \item Crop de deux fen\^etres centr\'ees sur les poignets ($\to$ deux
        sous-images).
  \item \textbf{MediaPipe Hands} sur chaque crop : 21 landmarks $(x, y)$
        par main d\'etect\'ee.
  \item \textbf{MLP} (42 $\to$ 30 $\to$ 23 logits) sur les 42 coordonn\'ees
        \og main principale \fg{} : $\arg\max$ + softmax $\to$ \texttt{letter}
        + \texttt{confidence}.
  \item R\'eponse JSON unique :
\begin{lstlisting}
{
  "skeleton":   [[x,y,vis], ... 9 pts],
  "hands":      [[[x,y], ... 21 pts], ...],
  "letter":     "A",
  "confidence": 0.84,
  "sign":       null,
  "sign_confidence": null
}
\end{lstlisting}
\end{enumerate}

\subsection*{C\^ot\'e client}

Le client dessine en permanence skeleton + hands. Pour la lettre :

\begin{itemize}[leftmargin=*, itemsep=2pt]
  \item \textbf{Seuil de confiance} : on ignore la pr\'ediction si
        \texttt{confidence < 0.6}.
  \item \textbf{Stabilit\'e 10 frames} : on n'ajoute la lettre au mot que
        si la m\^eme valeur sort 10 frames de suite. \'Evite que le bruit
        entre deux signes pollue le mot final.
  \item \textbf{Aucun seuil g\'eom\'etrique} : la ligne du seuil
        (canvas, slider) est masqu\'ee en mode statique --- l'inf\'erence
        tourne sans interruption.
\end{itemize}

\section{Mode dynamique --- GRU Ocarina + d\'eclencheur de seuil}

\textbf{Objectif} : reconna\^itre les signes \emph{en mouvement} (J, P, X, Z,
mots complets) qui n\'ecessitent une trajectoire temporelle.

\subsection*{Flux serveur}

L'envoi reste continu (m\^eme rythme, m\^eme protocole). Les \'etapes 1--4
ci-dessus sont identiques. La diff\'erence vient apr\`es l'\'etape 4 :

\begin{enumerate}[leftmargin=*, itemsep=2pt, start=5]
  \item Pour chaque session Socket.IO, le serveur maintient une
        \textbf{m\'emoire roulante} : un \texttt{deque(maxlen=45)} de
        vecteurs 42-float.
  \item \textbf{Normalisation des landmarks} avant push dans le buffer :
        \begin{itemize}[leftmargin=*, itemsep=0pt]
          \item rescale depuis la r\'esolution runtime vers l'espace
                pixel d'entra\^inement (640$\times$480),
          \item recentre tous les points sur le poignet (point~0) $\to$
                \textbf{invariance en translation}.
        \end{itemize}
  \item Toutes les \texttt{SIGN\_EVERY\_N = 5} frames, si le buffer est
        plein (45 frames disponibles), on appelle le \textbf{GRU Ocarina}
        sur la s\'equence compl\`ete.
\end{enumerate}

\paragraph{Architecture GRU Ocarina.} Embedding lin\'eaire $42 \to 256
\to 64$ avec ReLU + Dropout 0{,}2, puis un GRU \`a 1 couche cach\'ee
($\text{hidden}=64$, \texttt{batch\_first=True}), puis un Linear
$64 \to N_\text{classes}$. On prend uniquement la sortie du dernier
pas de temps. La pr\'ediction sort sous forme \texttt{sign} +
\texttt{sign\_confidence}.

R\'eponse JSON enrichie :

\begin{lstlisting}
{
  "skeleton":   [...],
  "hands":      [...],
  "letter":     "A",        // toujours calcul\'e (gratuit)
  "confidence": 0.65,
  "sign":       "Bonjour",  // <- nouveau, GRU Ocarina
  "sign_confidence": 0.82
}
\end{lstlisting}

\subsection*{Le d\'eclencheur de seuil c\^ot\'e client}

C'est ici que le mode dynamique se distingue r\'eellement : le client
ne fait pas confiance \`a chaque \texttt{sign} re\c{c}u, il attend la
\textbf{fin du geste} pour valider.

\begin{enumerate}[leftmargin=*, itemsep=3pt]
  \item Le navigateur charge \textbf{MediaPipe Hands en local} (CDN,
        modelComplexity=0, $\sim$12 MB cach\'es par le navigateur).
        Un RAF loop pompe la frame webcam locale et r\'ecup\`ere les
        landmarks \emph{c\^ot\'e client} ind\'ependamment du serveur.
  \item Une \textbf{ligne horizontale} est dessin\'ee sur un canvas
        au-dessus de la vid\'eo, \`a la hauteur \texttt{thresholdRatio}
        ($\in [0{,}10\ ;\ 0{,}90]$, r\'eglable via un slider, persist\'e
        en \texttt{localStorage}).
  \item \`A chaque tick RAF, on compare le $y$ minimum des landmarks
        (point le plus haut) au seuil :
        \begin{itemize}[leftmargin=*, itemsep=0pt]
          \item \textbf{au-dessus} $\to$ ligne verte (\og signe en
                cours \fg{}),
          \item \textbf{en-dessous} pendant plus de 500~ms continus
                $\to$ ligne rouge (\og en attente \fg{}).
        \end{itemize}
  \item \`A chaque pr\'ediction GRU re\c{c}ue du serveur avec
        \texttt{sign\_confidence} $\geq 0{,}7$, on \textbf{stocke}
        la pr\'ediction dans \texttt{pendingSign} (en gardant le sign
        de plus haute confidence si plusieurs arrivent).
  \item Au moment de la \textbf{transition au-dessus $\to$ en-dessous}
        (d\'ebounc\'ee 500~ms) : le client consid\`ere le geste termin\'e
        et \textbf{commit} \texttt{pendingSign} dans le mot.
\end{enumerate}

\paragraph{Pas de stabilit\'e 10 frames.} Contrairement au MLP, le
GRU travaille d\'ej\`a sur 45 frames temporellement coh\'erentes. Il
sort une pr\'ediction d\'ej\`a stable par construction. Imposer une
stabilit\'e suppl\'ementaire reviendrait \`a sommer du bruit.

\paragraph{Le squelette reste affich\'e en permanence.} M\^eme quand la
main est sous le seuil, le serveur continue de renvoyer
\texttt{skeleton} + \texttt{hands} et le client les redessine. Seul le
\emph{commit du sign dans le mot} est gat\'e par le seuil.

\section{Synth\`ese visuelle des deux flows}

\subsection*{Statique}

\begin{center}
\begin{tikzpicture}[
  font=\footnotesize, node distance=0.35cm,
  step/.style={draw, rounded corners=2pt, minimum height=7mm, minimum width=22mm, align=center, fill=bg},
  arr/.style={-{Stealth[length=1.8mm]}, line width=0.3mm},
]
  \node[step] (a) {frame webcam};
  \node[step, right=of a] (b) {Vision + MediaPipe};
  \node[step, right=of b] (c) {MLP};
  \node[step, right=of c] (d) {conf $\geq 0{,}6$ ?};
  \node[step, right=of d] (e) {10 frames \\ identiques ?};
  \node[step, right=of e, fill=good!20, draw=good] (f) {+1 lettre \\ au mot};
  \draw[arr] (a) -- (b); \draw[arr] (b) -- (c); \draw[arr] (c) -- (d);
  \draw[arr] (d) -- (e); \draw[arr] (e) -- (f);
\end{tikzpicture}
\end{center}

\subsection*{Dynamique}

\begin{center}
\begin{tikzpicture}[
  font=\footnotesize, node distance=0.35cm,
  step/.style={draw, rounded corners=2pt, minimum height=7mm, minimum width=22mm, align=center, fill=bg},
  arr/.style={-{Stealth[length=1.8mm]}, line width=0.3mm},
]
  \node[step] (a) {frame webcam};
  \node[step, right=of a] (b) {Vision + MediaPipe};
  \node[step, right=of b] (c) {push dans buffer \\ 45 frames};
  \node[step, right=of c] (d) {GRU Ocarina \\ (tous les 5)};
  \node[step, below=0.6cm of c, fill=accent2!20, draw=accent2!80!black] (cli1) {MediaPipe \\ client-side};
  \node[step, right=of cli1, fill=accent2!20, draw=accent2!80!black] (cli2) {$y_\text{main} < $ seuil ?};
  \node[step, right=of cli2, fill=accent2!20, draw=accent2!80!black] (cli3) {transition \\ haut $\to$ bas};
  \node[step, right=of cli3, fill=good!20, draw=good] (cli4) {+1 signe \\ au mot};
  \draw[arr] (a) -- (b); \draw[arr] (b) -- (c); \draw[arr] (c) -- (d);
  \draw[arr] (d) -- ++(2.2,0) |- node[pos=0.25, right]{\texttt{pendingSign}} (cli3);
  \draw[arr] (cli1) -- (cli2); \draw[arr] (cli2) -- (cli3); \draw[arr] (cli3) -- (cli4);
\end{tikzpicture}
\end{center}

\section{R\'ecapitulatif des constantes}

\begin{center}
\renewcommand{\arraystretch}{1.2}
\begin{tabular}{lll}
\hline
\textbf{Constante} & \textbf{Valeur} & \textbf{Effet} \\
\hline
\texttt{SEQ\_LEN}                 & 45    & taille du buffer temporel GRU \\
\texttt{SIGN\_EVERY\_N}           & 5     & inf\'erence GRU 1 frame sur 5 \\
\texttt{TRAIN\_W}, \texttt{TRAIN\_H} & 640, 480 & espace pixel canonique de normalisation \\
\texttt{MIN\_LETTER\_CONFIDENCE}  & 0{,}6 & seuil de confiance MLP \\
\texttt{STABLE\_FRAMES\_TO\_COMMIT}& 10   & stabilit\'e MLP avant commit \\
\texttt{MIN\_SIGN\_CONFIDENCE}    & 0{,}7 & seuil de confiance GRU \\
\texttt{ABSENCE\_GRACE\_MS}       & 500 ms & d\'ebounce de la descente sous le seuil \\
\texttt{thresholdRatio}           & 0{,}40 & hauteur du seuil (slider 10--90~\%) \\
\hline
\end{tabular}
\end{center}

\vspace{0.6em}
\noindent\textbf{Code source :} \url{https://github.com/nmqx/hearmyhands}\\
\textbf{Live :} \url{https://hearmyhands.asia}

\end{document}
